\documentclass{article}
\usepackage[utf8]{inputenc}
\usepackage{amsmath}
\usepackage{geometry}
\usepackage{float}
\usepackage{booktabs}
\geometry{a4paper, margin=1in}

\title{Homework 6 - Part 2c: Final Load and Center of Pressure Analysis}
\author{}
\date{}

\begin{document}
\maketitle

\section*{Part 2c: Summary and Findings}
Table \ref{tab:final_results} summarizes the computed load ($F$) and the normalized center of pressure ($x_p/L$) for the three slider configurations modeled in this assignment.

\begin{table}[H]
    \centering
    \renewcommand{\arraystretch}{1.3}
    \begin{tabular}{lcc}
        \toprule
        \textbf{Configuration Case} & \textbf{Load $F$ (N)} & \textbf{Center of Pressure $x_p/L$} \\
        \midrule
        \textbf{1a:} Incompressible 2D Taper-Flat & $0.2067$ & $0.5387$ \\
        \textbf{1b:} Incompressible 1D Pure Wedge & $8.1541$ & $0.5391$ \\
        \textbf{2a:} Compressible 2D Taper-Flat   & $0.2783$ & $0.5352$ \\
        \bottomrule
    \end{tabular}
    \caption{Comparison of the lifting load and normalized center of pressure across different fluid and geometric assumptions.}
    \label{tab:final_results}
\end{table}

\subsection*{Discussion of Findings}
Based on the results in Table \ref{tab:final_results}, several key physical insights can be drawn regarding the modeling of air bearings:

\begin{itemize}
    \item \textbf{The Impact of Side Leakage (1D vs. 2D):} 
    The 1D wedge model (Case 1b) predicts a load of $8.1541$ N, which is nearly 40 times larger than the realistic 2D model (Case 1a). By artificially restricting the flow to a single dimension (Neumann boundaries), the fluid cannot escape laterally. This creates an intense, artificial pressure ridge, grossly overestimating the slider's load-carrying capacity. This demonstrates why 2D modeling is strictly necessary for finite-width sliders.
    
    \item \textbf{The Effect of Compressibility (Case 1a vs. Case 2a):} 
    Introducing the compressible Reynolds equation (Case 2a) results in a load of $0.2783$ N, an increase of approximately 34\% over the incompressible baseline (Case 1a). As the air is squeezed into the narrow trailing gap, its absolute pressure rises, which proportionally increases the fluid density. This localized densification ``stiffens'' the air film, allowing it to generate more lift under the exact same geometric and kinematic conditions.
    
    \item \textbf{Stability of the Center of Pressure:} 
    Despite the massive variations in predicted load magnitudes, the normalized center of pressure ($x_p/L$) remains highly stable across all three models, situated at approximately $53.5\%$ to $53.9\%$ of the slider's length. The force consistently acts slightly downstream of the geometric center, which is characteristic of the wedge-like geometry squeezing fluid toward the trailing edge. 
\end{itemize}

\end{document}